\section{第3章の解答}\label{解答:第3章}
\paragraph{3.1}
$a=b$ なら両辺とも $0$ である.\,$a<b$ とすると,\,$f(x)=\sin x$ は $[a,b]$ で連続,\,$(a,b)$ で微分可能である.平均値の定理により,ある $c\in(a,b)$ について
\[
 \sin b-\sin a=(b-a)\cos c
\]
が成り立つ.\,$|\cos c|\le1$ だから,
\[
 |\sin b-\sin a|=|b-a|\,|\cos c|\le|b-a|.
\]
$a>b$ の場合は $a,b$ を入れ替えて同じ議論を使えばよい.

\paragraph{3.2}
導関数は
\[
 f'(x)=3x^2-6x=3x(x-2),\qquad f''(x)=6x-6
\]
である.したがって,\,$x<0$ で増加,\,$0<x<2$ で減少,\,$x>2$ で増加する.\,$x=0$ で極大値 $0$,\,$x=2$ で極小値 $-4$ をとる.また,\,$x<1$ で上に凸,\,$x>1$ で下に凸なので,変曲点は $(1,f(1))=(1,-2)$ である.

閉区間 $[-1,3]$ では,極値をとる点と両端を比較する.
\[
\begin{array}{c|rrrr}
 x&-1&0&2&3\\ \hline
 f(x)&-4&0&-4&0
\end{array}
\]
したがって最大値は $0$ で,\,$x=0,3$ のときにとる.最小値は $-4$ で,\,$x=-1,2$ のときにとる.

\paragraph{3.3}
$h(x)=\e^x-1-x$ とおくと,\,$h'(x)=\e^x-1$ である.\,$x<0$ では負,\,$x>0$ では正なので,\,$h$ は $x=0$ で最小値 $h(0)=0$ をとる.よって $\e^x-1-x\ge0$,すなわち $\e^x\ge1+x$ である.\,$0$ の両側で狭義に増減しているから,等号は $x=0$ に限る.

方程式については $p(x)=\log x+x-2$ とおく.\,$p$ は $x>0$ で連続かつ微分可能で,
\[
 p'(x)=\frac1x+1>0.
\]
したがって狭義に増加し,\,$0$ をとる点は高々1つである.一方,
\[
 p(1)=-1<0,\qquad p(2)=\log2>0
\]
だから,中間値の定理により $(1,2)$ に少なくとも1つの解がある.以上より,\,$x>0$ で解はただ1つであり,その解は $(1,2)$ にある.

\paragraph{3.4}
(1) $f'(x)=3x^2+1>0$ なので $f$ は狭義に増加する.また,\,$f(x)$ は $x\to\pm\infty$ のときそれぞれ $\pm\infty$ に向かうため,実数全体を値域とする逆関数をもつ.\,$f(0)=0$,\,$f'(0)=1$ から,
\[
 g'(0)=\frac1{f'(0)}=1.
\]
逆関数そのものを解の公式で表す必要はない.

(2) $t=2$ で $(x,y)=(4,8)$ であり,
\[
 \frac{\dd x}{\dd t}=2t,\qquad
 \frac{\dd y}{\dd t}=3t^2,\qquad
 \left.\frac{\dd y}{\dd x}\right|_{t=2}
 =\frac{12}{4}=3.
\]
よって接線は $y-8=3(x-4)$,すなわち $y=3x-4$ である.

(3) $y$ を $x$ の関数として両辺を微分すると,
\[
 2x+y+xy'+2yy'=0.
\]
$y'$ を含む項を集めて,
\[
 (x+2y)y'=-(2x+y),\qquad
 y'=-\frac{2x+y}{x+2y}
\]
を得る.点 $(1,2)$ では分母は $5\ne0$ で,傾きは $-\frac45$ となる.したがって接線は
\[
 y-2=-\frac45(x-1),\qquad 4x+5y=14
\]
である.元の関係式に点を代入して $1+2+4=7$ となることも確認できる.

\paragraph{3.5}
積の微分と$y'=3y$から
\begin{align*}
 (\e^{-3x}y)'&=-3\e^{-3x}y+\e^{-3x}y'\\
 &=\e^{-3x}(y'-3y)\\
 &=0
\end{align*}
である.したがって$\e^{-3x}y$は定数であり,\,$x=0$を代入するとその値は$2$となる.よって$y=2\e^{3x}$である.これを微分すれば$y'=6\e^{3x}=3y$となり,初期条件も満たす.
\paragraph{3.6}
$f(x)=x^2-3$,\,$f'(x)=2x$だから,
\begin{align*}
 x_{n+1}&=x_n-\frac{x_n^2-3}{2x_n}\\
 &=\frac12\left(x_n+\frac3{x_n}\right)
\end{align*}
である.\,$x_n>\sqrt3$なら
\[
 x_{n+1}-\sqrt3=\frac{(x_n-\sqrt3)^2}{2x_n}>0,
 \qquad x_{n+1}-x_n=\frac{3-x_n^2}{2x_n}<0
\]
となる.\,$x_1=2>\sqrt3$なので,帰納的に全項が$\sqrt3$より大きく,単調減少である.下に有界な単調減少数列なので,数列は収束する.その極限を$L$とすると$L\ge\sqrt3$であり,漸化式の両辺の極限をとって
\[
 L=\frac12\left(L+\frac3L\right)
\]
となる.これより$L^2=3$であり,\,$L>0$だから$L=\sqrt3$である.
